% !PRINTED 33
\weberpage{33}
%<<CONT
multiple of $a$. The element 1 has degree 1. All cyclic groups are commutative.

Now for every Abelian group $S$ there holds the fundamental theorem that a system of elements $A, B, C \ldots$ of degrees $a, b, c \ldots$ can always be selected in such a way that in the form
\[
\Theta = A^\alpha B^\beta C^\gamma \ldots
\]
every element $\Theta$ of $S$, and each only once, can be represented, when $\alpha$ runs through a complete system of residues to the modulus $a$, $\beta$ a complete system of residues to the modulus $b$, $\gamma$ a complete system of residues to the modulus $c$, etc.

A system of elements that is suited to such a representation is called a \gesperrt{basis of the group}, and we can accordingly state the theorem to be proved thus:

\gesperrt{I. Every Abelian group can be represented by a basis.}

The proof of the theorem rests on the following sequence of inferences:

\gesperrt{1. If $n$ is the degree of the group $S$, if $A_1, A_2 \ldots A_n$ are its elements and $a_1, a_2 \ldots a_n$ their degrees, and if further $\alpha_1, \alpha_2 \ldots \alpha_n$ run through complete systems of residues to the moduli $a_1, a_2 \ldots a_n$, then in the form}
\weq{1}{\Theta = A_1^{\alpha_1} A_2^{\alpha_2} \ldots A_n^{\alpha_n}}
\gesperrt{every element of $S$ is represented, and each equally often.}

That one obtains every element at all in the form (1) is seen immediately; for one obtains, e.\,g., $A_1$ by setting $\alpha_1 = 1, \alpha_2 = 0 \ldots \alpha_n = 0$.

It remains to be proved that one obtains every element equally often.

But if
\weq{2}{1 = A_1^{h_1} A_2^{h_2} \ldots A_n^{h_n}}
is a representation of the element 1, then $\Theta$ does not change when $\alpha_1, \alpha_2 \ldots \alpha_n$ in (1) are replaced by $\alpha_1 + h_1, \alpha_2 + h_2 \ldots \alpha_n + h_n$. Thus every element $\Theta$ is represented by (1) at least as often as the element 1 is by (2).

If, on the other hand, the two sequences of exponents $\alpha_1, \alpha_2 \ldots \alpha_n$ and $\alpha_1', \alpha_2' \ldots \alpha_n'$ give two representations of the element $\Theta$, then one has
%CONT>>
